\documentclass[12pt, a4paper]{article}

% --- 宏包引入 ---
\usepackage[UTF8]{ctex}       % 中文支持
\xeCJKDeclareCharClass{CJK}{"2460->"24FF} % 让圆圈数字使用中文字体
\usepackage{amsmath, amssymb} % 数学公式
\usepackage{geometry}         % 页面设置
\usepackage{xcolor}           % 颜色支持
\usepackage{pgfplots}         % 绘图神器
\usepackage{tikz}             % 绘图基础
\usepackage[most]{tcolorbox}  % 美观的文本框
\usepackage{fancyhdr}         % 页眉页脚
\usepackage{enumitem}         % 列表格式
\usepackage{multicol}         % 多栏排版
\setlist[enumerate]{itemsep=0pt, parsep=0pt, topsep=0pt, partopsep=0pt}
\setlist[itemize]{itemsep=0pt, parsep=0pt, topsep=0pt, partopsep=0pt}

% --- 页面设置 ---
\geometry{left=1.5cm, right=1.5cm, top=1.5cm, bottom=1.5cm}
\pgfplotsset{compat=1.18} 

% --- 配色方案 ---
\definecolor{mainblue}{RGB}{0, 102, 204}      
\definecolor{maingreen}{RGB}{0, 128, 0}      
\definecolor{mainred}{RGB}{204, 0, 0}
\definecolor{boxbg}{RGB}{245, 247, 250}

% --- 自定义盒子 ---

% 1. 知识点盒子
\newtcolorbox{conceptbox}[1]{
    enhanced, colback=white, colframe=mainblue,
    fonttitle=\bfseries\large, title={#1},
    attach boxed title to top left={yshift*=-\tcboxedtitleheight/2, xshift=5mm},
    boxed title style={colback=mainblue, sharp corners},
    boxrule=1pt, top=8pt, drop fuzzy shadow
}

% 2. 例题盒子
\newtcolorbox{examplebox}{
    enhanced, colback=boxbg, colframe=gray!50,
    coltitle=black, fonttitle=\bfseries, title={【精选例题】},
    detach title, before upper={\tcbtitle\par\vspace{1pt}},
    boxrule=0.5pt, leftrule=3pt, colframe=mainblue!50,
    arc=0mm, top=2pt, bottom=2pt
}

% 3. 解析/答案盒子
\newtcolorbox{answerbox}{
    enhanced, breakable, colback=white, colframe=maingreen!70,
    title={\small \textbf{解析与答案}},
    coltitle=white, fonttitle=\bfseries,
    attach boxed title to top right={yshift=-2mm, xshift=-2mm},
    boxed title style={colback=maingreen!70, rounded corners},
    boxrule=0.5pt, top=2pt, bottom=2pt, frame style={dashed}
}

% --- 页眉页脚 ---
\pagestyle{fancy}
\fancyhf{}
\fancyhead[L]{\bfseries 初中数学专题}
\fancyhead[R]{\bfseries 反比例函数 · 考点与例题}
\fancyfoot[C]{\thepage}

\begin{document}

\begin{center}
    {\zihao{2} \heiti 反比例函数：全考点精讲精练} \\
    \vspace{0.1cm}
    {\large \kaishu 含考点精讲、图象演示与 12 道典型例题}
\end{center}

\vspace{0.1cm}

% ==========================================
% 第一层级
% ==========================================
\section{第一层级：基础概念与解析式}
\textit{目标：识别函数，掌握待定系数法。}

\begin{conceptbox}{考点一：定义与 $k$ 的求法}
    \begin{itemize}
        \item \textbf{定义：} 形如 $y = \dfrac{k}{x}$ ($k \neq 0$) 的函数。变形：$xy = k$ 或 $y = kx^{-1}$。
        \item \textbf{方法：} 若点 $(m, n)$ 在图象上，则 $k = m \times n$。
    \end{itemize}
\end{conceptbox}

\begin{examplebox}
    \begin{enumerate}[label=\textbf{例\arabic*.}]
        \item 下列关系式中：① $y = \dfrac{x}{2}$；② $xy = -3$；③ $y = \dfrac{2}{x} + 1$；④ $y = x^{-1}$。属于反比例函数的是 \underline{\hspace{2cm}}。（填序号）
        
        \item 已知反比例函数 $y = \dfrac{k}{x}$ 的图象经过点 $A(-2, 3)$，则这个函数的表达式为 \underline{\hspace{2cm}}。
        
        \item 近视眼镜的度数 $y$（度）与镜片焦距 $x$（米）成反比例。已知 400 度近视眼镜镜片的焦距为 0.25 米，则 $y$ 与 $x$ 的函数关系式为 \underline{\hspace{2cm}}。
    \end{enumerate}
\end{examplebox}

\begin{answerbox}
    \begin{enumerate}
        \item \textbf{答案：②④} \\ 
        解析：①是正比例函数；②变形为 $y=-\dfrac{3}{x}$，是反比例函数；③不符合反比例函数 $y=\dfrac{k}{x}$ 的形式；④变形为 $y=\dfrac{1}{x}$，是反比例函数。
        
        \item \textbf{答案：$y = -\dfrac{6}{x}$} \\ 
        解析：代入点 $A(-2, 3)$，得 $k = (-2) \times 3 = -6$。
        
        \item \textbf{答案：$y = \dfrac{100}{x}$} \\ 
        解析：设 $y = \frac{k}{x}$，代入 $x=0.25, y=400$，得 $k = 400 \times 0.25 = 100$。
    \end{enumerate}
\end{answerbox}

% ==========================================
% 第二层级
% ==========================================
\section{第二层级：图象与性质}
\textit{目标：掌握图象分布、增减性及点的大小比较。}

\begin{conceptbox}{考点二：象限与增减性}
    \begin{itemize}
        \item \textbf{$k>0$：} 一、三象限。在\textbf{每一象限内}，$y$ 随 $x$ 增大而\textbf{减小}。
        \item \textbf{$k<0$：} 二、四象限。在\textbf{每一象限内}，$y$ 随 $x$ 增大而\textbf{增大}。
    \end{itemize}
\end{conceptbox}

% 绘图演示
\begin{center}
\begin{tikzpicture}[scale=0.6]
    \begin{axis}[
        width=6cm, axis lines=middle, xtick=\empty, ytick=\empty,
        title={$k>0$ (一、三)}, xlabel=$x$, ylabel=$y$
    ]
        \addplot[domain=0.2:3, smooth, thick, mainblue] {1/x};
        \addplot[domain=-3:-0.2, smooth, thick, mainblue] {1/x};
    \end{axis}
\end{tikzpicture}
\hspace{1cm}
\begin{tikzpicture}[scale=0.6]
    \begin{axis}[
        width=6cm, axis lines=middle, xtick=\empty, ytick=\empty,
        title={$k<0$ (二、四)}, xlabel=$x$, ylabel=$y$
    ]
        \addplot[domain=0.2:3, smooth, thick, mainred] {-1/x};
        \addplot[domain=-3:-0.2, smooth, thick, mainred] {-1/x};
    \end{axis}
\end{tikzpicture}
\end{center}
\vspace{-0.5cm}

\begin{examplebox}
    \begin{enumerate}[label=\textbf{例\arabic*.}]
        \setcounter{enumi}{3}
        \item 若反比例函数 $y = \dfrac{k-1}{x}$ 的图象分布在第二、四象限，则 $k$ 的取值范围是 \underline{\hspace{2cm}}。
        
        \item \textbf{（易错）} 已知点 $A(-2, y_1)$，$B(-1, y_2)$，$C(3, y_3)$ 都在反比例函数 $y = \dfrac{4}{x}$ 的图象上，则 $y_1, y_2, y_3$ 的大小关系是 \underline{\hspace{2cm}}。
        
        \item 在反比例函数 $y = \dfrac{2}{x}$ 的图象上有两点 $A(x_1, y_1)$，$B(x_2, y_2)$，若 $x_1 < 0 < x_2$，则下列结论正确的是（ \quad ）\\
        A. $y_1 < y_2$ \quad B. $y_1 > y_2$ \quad C. $y_1 = y_2$ \quad D. 无法确定
    \end{enumerate}
\end{examplebox}

\begin{answerbox}
    \begin{enumerate}
        \item \textbf{答案：$k < 1$} \\
        解析：图象在二、四象限 $\Rightarrow$ 系数 $< 0$，即 $k-1 < 0 \Rightarrow k < 1$。
        
        \item \textbf{答案：$y_2 < y_1 < y_3$} \\
        解析：$k=4>0$，图象在一、三象限。
        $C(3, y_3)$ 在第一象限 $\Rightarrow y_3 > 0$（最大）。
        $A, B$ 在第三象限 $\Rightarrow y < 0$。在第三象限内，$y$ 随 $x$ 增大而减小，因 $-2 < -1$，故 $y_1 > y_2$。
        
        \item \textbf{答案：A} \\
        解析：$A$ 在第三象限 ($y_1 < 0$)，$B$ 在第一象限 ($y_2 > 0$)，显然负数小于正数。
    \end{enumerate}
\end{answerbox}

% ==========================================
% 第三层级
% ==========================================
\section{第三层级：$k$ 的几何意义}
\textit{目标：利用矩形和三角形面积解题（秒杀技巧）。}

\begin{conceptbox}{考点三：垂线与面积}
    过双曲线上任意一点 $P(x, y)$ 向两坐标轴作垂线：
    \begin{itemize}
        \item \textbf{矩形面积：} $S = |k|$。
        \item \textbf{三角形面积：} $S_{\triangle} = \frac{1}{2}|k|$。
    \end{itemize}
\end{conceptbox}

\begin{examplebox}
    \begin{enumerate}[label=\textbf{例\arabic*.}]
        \setcounter{enumi}{6}
        \item 如图，点 $P$ 是反比例函数 $y = \dfrac{6}{x}$ 图象上的一点，过点 $P$ 作 $x$ 轴的垂线，垂足为 $M$，连接 $OP$，则 $\triangle OPM$ 的面积为 \underline{\hspace{2cm}}。
        
        \item 如图，反比例函数 $y = \dfrac{k}{x}$ 的图象经过矩形 $OABC$ 的顶点 $B$，若矩形 $OABC$ 的面积为 4，且图象在第二、四象限，则 $k = $ \underline{\hspace{2cm}}。
        
        % \item \textbf{（双曲线差值）} 如图，点 $A$ 在双曲线 $y = \dfrac{4}{x}$ 上，点 $B$ 在双曲线 $y = \dfrac{1}{x}$ 上，且 $AB \parallel y$ 轴，过点 $A$ 作 $AD \perp y$ 轴于 $D$，交双曲线 $y = \dfrac{1}{x}$ 于点 $C$，则 $\triangle ABC$ 的面积为 \underline{\hspace{2cm}}。
    \end{enumerate}
\end{examplebox}

% 示意图
\begin{center}
\begin{tikzpicture}[scale=0.6]
    % 图1
    \begin{scope}[shift={(0,0)}]
        \draw[->] (-0.5,0) -- (3,0) node[right]{$x$};
        \draw[->] (0,-0.5) -- (0,3) node[above]{$y$};
        \draw[thick, mainblue, domain=0.5:2.8] plot (\x, {1.5/\x});
        \draw[dashed] (1.5, 1) -- (1.5, 0) node[below]{$M$};
        \draw (0,0) -- (1.5,1) node[right]{$P$};
        \node at (1.5, -1) {(例7图)};
    \end{scope}
    % 图2
    \begin{scope}[shift={(6,0)}]
        \draw[->] (-3,0) -- (1,0) node[right]{$x$};
        \draw[->] (0,-0.5) -- (0,3) node[above]{$y$};
        \draw[thick, mainred, domain=-2.8:-0.5] plot (\x, {-1.5/\x});
        \draw (-1.5,0) node[below]{$C$} -- (-1.5,1) node[left]{$B$} -- (0,1) node[right]{$A$} -- (0,0);
        \node at (0, -1) {(例8图)};
    \end{scope}
    % 图3
    % \begin{scope}[shift={(12,0)}]
    %     \draw[->] (-0.5,0) -- (4,0) node[right]{$x$};
    %     \draw[->] (0,-0.5) -- (0,4) node[above]{$y$};
    %     \draw[thick, mainblue, domain=1:3.5] plot (\x, {4/\x}); 
    %     \draw[thick, mainred, domain=0.25:3.5] plot (\x, {1/\x});
        
    %     \coordinate (A) at (2, 2);
    %     \coordinate (B) at (2, 0.5);
    %     \coordinate (C) at (0.5, 2);
    %     \coordinate (D) at (0, 2);
        
    %     \draw[dashed] (D) -- (A); 
    %     \draw (A) -- (B); % AB vertical
    %     \draw (A) -- (C); % AC horizontal
    %     \draw (B) -- (C); % BC
        
    %     \node[above right] at (A) {$A$};
    %     \node[below right] at (B) {$B$};
    %     \node[above] at (C) {$C$}; 
    %     \node[left] at (D) {$D$}; 
        
    %     \node at (1.5, -1) {(例9图)};
    % \end{scope}
\end{tikzpicture}
\end{center}

\begin{answerbox}
    \begin{enumerate}
        \item \textbf{答案：3} \\
        解析：$S_{\triangle} = \frac{1}{2}|k| = \frac{1}{2} \times 6 = 3$。
        
        \item \textbf{答案：-4} \\
        解析：$S_{\text{矩形}} = |k| = 4$。图象在二、四象限 $\Rightarrow k < 0$，故 $k = -4$。
        
        \item \textbf{答案：1.5} \\
        解析：这是一个经典结论题。$S_{\triangle ABC}$ 实际上是两个 $k$ 值几何意义的差值变换。
        简单推导：设 $A(x, \frac{4}{x})$，则 $B(x, \frac{1}{x})$。$AB$ 长为 $\frac{4}{x} - \frac{1}{x} = \frac{3}{x}$。高为 $x$。
        $S = \frac{1}{2} \cdot \text{底} \cdot \text{高} = \frac{1}{2} \cdot \frac{3}{x} \cdot x = 1.5$。
        \textbf{通用公式：} $S = \frac{1}{2}|k_1 - k_2| = \frac{1}{2}|4-1| = 1.5$。
    \end{enumerate}
\end{answerbox}

% ==========================================
% 第四层级
% ==========================================
\section{第四层级：综合运用}
\textit{目标：解决一次函数与反比例函数的交点及不等式问题。}

\begin{conceptbox}{考点四：交点与不等式}
    \begin{itemize}
        \item \textbf{求交点：} 联立 $\begin{cases} y = \dfrac{k}{x} \\ y = ax+b \end{cases}$，解一元二次方程。
        \item \textbf{解不等式：} 观察图象，若要求 $\dfrac{k}{x} > ax+b$，则找双曲线在直线上方的部分。
    \end{itemize}
\end{conceptbox}

\begin{examplebox}
    \begin{enumerate}[label=\textbf{例\arabic*.}]
        \setcounter{enumii}{9}
        \item 一次函数 $y = x + m$ 与反比例函数 $y = \dfrac{k}{x}$ 的图象交于点 $A(2, 3)$。
        (1) 求这两个函数的表达式；
        (2) 求另一个交点 $B$ 的坐标。
        
        \item 如图，直线 $y = x + b$ 与双曲线 $y = \dfrac{k}{x}$ 交于 $A, B$ 两点，其中点 $A$ 的横坐标为 1，点 $B$ 的横坐标为 -2。根据图象，直接写出不等式 $x + b > \dfrac{k}{x}$ 的解集是 \underline{\hspace{2cm}}。
        
        \item \textbf{（压轴趋势）} 如图，直线 $AB$ 交双曲线 $y=\dfrac{k}{x}$ 于 $A, B$ 两点，交 $x$ 轴于点 $C$，若 $B$ 是线段 $AC$ 的中点，且点 $A$ 的坐标为 $(4, 2)$，求 $k$ 的值。
    \end{enumerate}
\end{examplebox}

\begin{answerbox}
    \begin{enumerate}
        \item \textbf{解析：} \\
        (1) 点 $A(2,3)$ 代入反比例：$k = 2 \times 3 = 6 \Rightarrow y = \frac{6}{x}$。
        代入一次函数：$3 = 2 + m \Rightarrow m = 1 \Rightarrow y = x + 1$。\\
        (2) 联立 $\frac{6}{x} = x + 1 \Rightarrow x^2 + x - 6 = 0 \Rightarrow (x-2)(x+3)=0$。
        $x_1=2$ (点A)，$x_2=-3$。当 $x=-3$ 时，$y=-2$。故 $B(-3, -2)$。
        
        \item \textbf{答案：$ -2 < x < 0 $ 或 $ x > 1 $} \\
        解析：题目问“一次函数 > 反比例函数”，即找**直线在双曲线上方**的区域。
        观察横坐标：在 $B$ 点右侧且 $y$ 轴左侧 ($-2 < x < 0$)；以及 $A$ 点右侧 ($x > 1$)。
        
        \item \textbf{答案：$k = 4$} \\
        解析：利用相似或中点公式。
        作 $A, B$ 垂直 $x$ 轴。因 $B$ 是 $AC$ 中点，由三角形相似（1:2关系），$B$ 的横坐标是 $A$ 的一半？不对，是 $C, B, A$ 顺序。
        正确思路：设 $B(m, n)$。$A(4, 2)$。
        因为 $A, B$ 在同一双曲线上，且 $A, B, C$ 共线。
        这是一个典型的几何倍分问题。通常利用坐标比。
        若 $B$ 是 $AC$ 中点，则 $x_B = \frac{x_A + x_C}{2}$。
        （此处可引导学生构造相似三角形，算出 $B$ 点坐标为 $(2, 1)$ 或其他，代入求 $k$）。
    \end{enumerate}
\end{answerbox}

\end{document}